\section{Self-Play and the Training Data}
\label{sec:selfplay}

Self-play runs in the Go engine, which is written in Go~\citep{donovan2015go}. This
section covers how a game is played, what each position is labelled with, and the binary
\emph{shard} format that replaces JSONL between the engine and the learner.

\subsection{Search during self-play}
\label{sec:selfplay-search}

Each move is chosen by Monte Carlo tree search with the PUCT selection rule of
\cref{eq:puct}~\citep{rosin2011multi,silver2017mastering}, using $c_{\mathrm{PUCT}}=1.1$
as in \citet{wu2019accelerating}. Two details differ from \katago{}. An unvisited child is
given the fixed value $-0.2$, whereas \katago{}'s first-play urgency is the parent's value
minus $0.2\sqrt{P_{\mathrm{explored}}}$. And the constant $1.03$, which \katago{} and
SAI~\citep{morandin2019sai} use as a softmax temperature on the root policy, is applied
in \gofer{} as a divisor on the root children's mean values. Root priors are mixed with Dirichlet noise ($\alpha=0.15$, weight $0.25$) on \emph{every}
self-play move. \katago{} instead turns noise and other exploration off on fast searches
to keep them as strong as possible~\citep[\S3.1]{wu2019accelerating}.

\paragraph{Playout cap randomization.}
Following \citet[\S3.1]{wu2019accelerating}, each move is a \emph{full} search (cap $N$)
with probability $p$, and otherwise a \emph{fast} search (cap $n<N$). The \gofer{}
defaults are $p=0.20$, $N=200$ and $n=50$; \katago{}'s main run used $p=0.25$ and
$(N,n)=(600,100)$, later raised to $(1000,200)$. The configurations in
\cref{sec:infra} change these values per deployment.

\paragraph{Forced playouts and target pruning (simplified).}
Full searches force a minimum number of visits to every root child $c$ that has been
expanded:
\begin{equation}
  n_{\mathrm{forced}}(c) = k + \left\lfloor \sqrt{P(c)\,(N+1)} \right\rfloor, \qquad k = 2,
  \label{eq:gofer-forced}
\end{equation}
where $P(c)$ is the noised prior. \katago{} uses
$n_{\mathrm{forced}}(c)=\bigl(kP(c)\sum_{c'}N(c')\bigr)^{1/2}$ and forces these visits by
setting the urgency of an under-visited child to infinity inside the search. \gofer{}
instead runs the forced playouts as a separate pass, right after the root is expanded and
before the main search.
The policy target differs more. \katago{} subtracts forced visits from each child only as
long as that child would not have overtaken the best child under PUCT. \gofer{} simply
removes children with fewer than two visits and renormalizes. This is cruder: it removes
single-visit noise but keeps any extra visits that forcing added to a child that stayed
bad. Both mechanisms predate v4 and are unchanged by it.

\paragraph{Two corrections to the search.}
Validating the v4 gate uncovered two defects in this search, both predating v4.
Selection used a child's mean without negating it, although node values are kept
from each node's own side to move, so the engine was choosing the move best for
its opponent; and two consecutive passes were not treated as the end of the
game, so a finished position was scored by the evaluator rather than by the
rules. \Cref{sec:pipeline-found} gives the evidence and the measured effect.
Everything described in this section is the corrected search.

\paragraph{Move selection and parallelism.}
For the first 16 plies the move is sampled in proportion to visit counts ($\tau=1$); after
that the most-visited move is played. The policy target is always the full visit
distribution, whatever move is played. Games run concurrently and share one batched
evaluator, so the neural network sees real batches. Each game's random generator is seeded
from the run seed and the game index only, so the output does not depend on how many games
run in parallel. It can still depend on the inference backend if two backends return
different network outputs; in our runs the two backends produced identical row counts
from the same seeds (\cref{sec:pipeline-eval}). Seeds are derived by mixing the run seed
with the game index rather than by adding them: consecutive seeds draw correlated
openings, which costs game diversity here and, in the arena, biased results outright
(\cref{sec:pipeline-found}).

\subsection{Labels and a common frame}
\label{sec:selfplay-labels}

Each recorded position gets the 8 binary input planes and 4 global features of
\cref{sec:learner-data}, plus these targets:
\begin{itemize}
  \item $\pi$: the root visit distribution over the $S^2+1$ moves (pass is the last entry);
  \item $z \in \{-1,0,+1\}$: the game result;
  \item $s$: the final area-score margin, with komi included;
  \item $o \in \{-1,0,+1\}^{S^2}$: the final owner of each point;
  \item $\pi_{\mathrm{opp}}$: the policy target of the \emph{next} ply, i.e.\ the
        opponent's reply, as in \citet{tian2016better} and \citet[\S3.4]{wu2019accelerating};
  \item a flag saying whether the move was a full search.
\end{itemize}
In v4, every target is expressed from the side to move, the same frame as the
``own/opponent'' input planes. Under area scoring every point belongs to one player or to
neither, so each row satisfies an exact identity:
\begin{equation}
  s \;=\; \sum_{i=1}^{S^2} o_i \;-\; \kappa\,\sigma, \qquad
  \sigma = \begin{cases} +1 & \text{Black to move}\\ -1 & \text{White to move}\end{cases}
  \label{eq:score-identity}
\end{equation}
Here $\kappa$ is the komi, which goes to White. The shard writer's test checks
\cref{eq:score-identity} on every row of real self-play output, and the fixture shard
shared with the learner's tests satisfies it with zero error.

\paragraph{Two defects in the v3 labels.}
Comparing the v3 JSONL fields with this frame found two errors. Neither raised an
exception, and both went unnoticed in v3.
\begin{enumerate}
  \item \textbf{Ownership frame.} v3 wrote ownership in the absolute frame (Black $=+1$)
        for every row, while value and the input planes are side-to-move. For
        White-to-move rows the target pointed the wrong way: the network had to recover
        the colour from the side-to-move plane just to learn a sign flip.
  \item \textbf{Opponent policy.} The field named \code{policy\_opp} held the policy of
        the \emph{previous} ply, not the opponent's reply. It is correlated with the
        right target, which is why nothing looked wrong. But the auxiliary loss it was
        meant to feed would have taught the network to predict the past.
\end{enumerate}
v4 fixes both in the engine. The shard format
stores the side-to-move ownership and the true next-ply policy. The JSONL writer keeps
its old field meanings for compatibility and adds \code{policy\_next} and
\code{score\_margin}; the learner corrects legacy rows when it loads them
(\cref{sec:learner-data}). $\pi_{\mathrm{opp}}$ is recorded only when the next ply was a
full search. After a fast search the visit distribution is too noisy to imitate, and the
row stores an all-zero vector, which the loss treats as missing.

\paragraph{Keeping fast-search rows.}
\katago{} records only full-search positions~\citep[\S3.1]{wu2019accelerating}, and v3 did
the same (\code{-full-only}). v4 shards keep every position and flag fast ones with
$\code{full\_search}=0$. The learner then masks the two policy losses to full-search rows
and trains value, ownership and score on every row (\cref{sec:learner-masking}). With
$p=0.2$ this gives about $5\times$ more rows per game for those heads, at no extra
self-play cost. The benefit is smaller than that ratio suggests. \citet{wu2019accelerating}
note that value training is limited by the number of \emph{games}, since there is one
noisy result per game. Extra positions from the same game share the same $z$ and are
strongly correlated. The gain we expect is therefore mostly in the per-point ownership
signal and in the variety of positions seen, not in independent value labels. The risk
is overfitting to game outcomes. The learner's game-level validation split
(\cref{sec:learner-data}) is what would reveal it. This choice has not been ablated for
strength (\cref{sec:limitations}).

\subsection{The shard format}
\label{sec:selfplay-shards}

When the output path ends in \code{.npz}, the engine writes a \emph{Gofer shard v1}. This
is a standard NumPy archive~\citep{harris2020array}: a ZIP file~\citep{pkware2022zip} of
\code{.npy} arrays~\citep{kern2007npy}, compressed with DEFLATE~\citep{deutsch1996deflate}.
There is one shard per self-play run, and every row in a shard has the same board size.
\Cref{tab:selfplay-shard} lists the arrays. Any \code{np.load} can read a shard, so the
learner, the replay index and ad-hoc analysis all use it without a custom parser. The
learner's JSONL packer produces the same layout, so the two producers are
interchangeable.

\begin{table}[htbp]
\centering
\small
\begin{tabular}{@{}llll@{}}
\toprule
Array & dtype & Shape & Content \\
\midrule
\code{spatial}      & uint8   & $[R,8,S,S]$ & binary input planes \\
\code{globals}      & float32 & $[R,4]$     & komi$/10$, move$/(S^2{+}1)$, is-black, is-white \\
\code{policy}       & float32 & $[R,S^2{+}1]$ & $\pi$, root visit distribution \\
\code{policy\_opp}  & float32 & $[R,S^2{+}1]$ & $\pi_{\mathrm{opp}}$, next-ply target (zeros if absent) \\
\code{value}        & float32 & $[R]$       & $z$, side to move \\
\code{score}        & float32 & $[R]$       & $s$, side to move, komi included \\
\code{ownership}    & int8    & $[R,S^2]$   & $o$, side to move \\
\code{full\_search} & uint8   & $[R]$       & 1 if the move was a full search \\
\code{game\_id}     & int32   & $[R]$       & game index within the shard \\
\code{move\_num}    & int16   & $[R]$       & ply number \\
\code{meta}         & uint8   & $[K]$       & UTF-8 JSON: format, version, board size, rows, games, \\
                    &         &             & komi, seed, generating model (ONNX hash), commit \\
\bottomrule
\end{tabular}
\caption{Arrays in a Gofer shard v1 with $R$ rows on an $S\times S$ board. Every target
uses the side-to-move frame.}
\label{tab:selfplay-shard}
\end{table}

Three properties matter for the rest of the pipeline:
\begin{itemize}
  \item \emph{Atomic writes.} The engine writes to a temporary file and renames it at the
        end, so a crash never leaves a truncated shard for the trainer to read.
  \item \emph{Immutability.} No later stage rewrites a shard, so the replay buffer is just a
        choice of which shards to read (\cref{sec:orch-replay}), and copying a run
        directory between machines is incremental.
  \item \emph{Provenance.} The \code{meta} array records which network generated the games
        (the first 8 bytes of its ONNX SHA-256, or \code{heuristic}). The orchestrator can
        read it without decompressing any data array.
\end{itemize}

\paragraph{Size.}
We measured size on the same four self-play games (seed 1, 20 playouts, 327 positions)
written in both formats. JSONL~\citep{bray2017json} took 779{,}565 bytes (2{,}384 bytes per
row), and the shard took 18{,}969 bytes (58 bytes per row), $41\times$ smaller.
Uncompressed, a shard row is 1{,}416 bytes, dominated by the 648 bytes of input planes and
the two 328-byte policy vectors. DEFLATE does the rest: the planes are binary and sparse,
and most policy entries are exactly zero. The learner's packer, run on a larger legacy file
(\cref{sec:eval}), gives a similar $38\times$. Loading a shard is one decompression per
array, with no per-number text parsing.
